\documentclass[11pt]{article}
\usepackage[T1]{fontenc}
\usepackage{lmodern,amsmath,amssymb,amsthm,mathtools}
\usepackage[margin=1in]{geometry}
\usepackage{microtype,booktabs,array,longtable}
\usepackage[numbers,sort&compress]{natbib}
\usepackage{xurl}
\usepackage[colorlinks=true,linkcolor=blue,citecolor=blue,urlcolor=blue]{hyperref}
\hypersetup{pdftitle={Sharp gaps for positive multilinear relaxations},
 pdfsubject={A counterexample to the Luedtke--Namazifar--Linderoth conjecture and Lean verification}}
\newtheorem{theorem}{Theorem}
\newtheorem{lemma}[theorem]{Lemma}
\newtheorem{proposition}[theorem]{Proposition}
\theoremstyle{remark}
\newtheorem{remark}[theorem]{Remark}
\DeclareMathOperator{\vex}{vex}
\DeclareMathOperator{\cav}{cav}
\DeclareMathOperator{\conv}{conv}
\newcommand{\E}{\mathbb E}
\newcommand{\Prb}{\mathbb P}
\newcommand{\R}{\mathbb R}
\newcommand{\ind}{\mathbf 1}
\newcommand{\laws}{\mathcal P}
\newcommand{\leanfile}[2]{\href{formal/Formal/#1.lean}{\texttt{#2}}}
\title{Sharp gaps for positive multilinear relaxations}
\author{}
\date{}
\begin{document}
\maketitle
\begin{abstract}
Luedtke, Namazifar, and Linderoth conjectured that, for positive-coefficient
multilinear functions on nonnegative boxes, the term-by-term relaxation gap
is within a universal constant of the convex-hull gap. We give a sparse
unit-coefficient family that disproves this conjecture. Its prescribed
point is strictly interior, its termwise gap is $L$, and its exact hull
gap is $s+(L-s)2^{-s}$ for an explicitly specified cutoff $s$.
A finite inequality and an attaining probability law prove the formula.
We also determine the sharp leading growth of the worst ratio:
$R(d)\sim\ln d/\ln\ln d$ in maximum degree and
$C(n)\sim\ln n/\ln\ln n$ in ambient dimension. The universal upper bound
uses one distribution for all terms and extends to the original termwise
relaxation on every finite nonnegative box. The accompanying standalone
Lean project verifies the disproof, exact hull-gap formula, and both
leading-constant-one limits. A claim-to-declaration guide and reproduction
instructions accompany the proofs.
\end{abstract}

\section{The conjecture and the gap being compared}
Let $B=\prod_{i=1}^n[\ell_i,u_i]$ be a finite box with
$0\le\ell_i\le u_i$. Consider a specified multilinear polynomial
\begin{equation}\label{eq:polynomial}
 f(x)=\sum_{e\in E}a_e m_e(x),\qquad m_e(x)=\prod_{i\in e}x_i,
 \qquad a_e\ge0,
\end{equation}
where $E$ is a set of distinct subsets of $[n]$. Zero coefficients can be
deleted; thus the positive-coefficient class is included. Empty supports
and singleton supports are allowed and contribute only affine terms.
For a continuous function $g$ on $B$, write $\vex_B g$ for its greatest
convex underestimator and $\cav_B g$ for its least concave overestimator.
Define
\begin{align}
 H_B(f;x)&=\cav_B f(x)-\vex_B f(x),\label{eq:hull-gap}\\
 T_B(f;x)&=\sum_{e\in E}
       [\cav_B(a_em_e)(x)-\vex_B(a_em_e)(x)].\label{eq:term-gap}
\end{align}
The first quantity is the vertical width of the convex hull of the graph
of $f$. The second gives each original monomial its exact graph hull
separately. Always $0\le H_B\le T_B$. We form $T_B/H_B$ only when $H_B>0$.

Conjecture~4.1 of \citet[p.~349]{Luedtke2012} proposes a constant upper
bound on this ratio, uniform over positive multilinear functions on
nonnegative boxes. The same statement is Conjecture~1 on printed page~22
of the linked author manuscript. The uniformity includes varying degree
and dimension. The construction below stays on unit cubes and uses
coefficient one for every included term.

The source distinguishes exact termwise envelopes from recursive
McCormick relaxations of higher-degree products on general boxes.
Our upper bounds concern \eqref{eq:term-gap}. They do not assert the same
bounds for every recursive formulation on a positive box. The mixed-sign
bilinear gap results of \citet{Boland2017} address a different coefficient
class. A pointwise width ratio also does not by itself give an
approximation ratio for a constrained optimization problem or an algorithm
for evaluating a convex envelope.

For $d\ge2$, let $R(d)$ be the supremum of $T_B/H_B$ over
\eqref{eq:polynomial} of maximum degree at most $d$, arbitrary finite
dimension, all such boxes, and all points with positive hull gap.
Let $C(n)$ be the analogous supremum over exactly $n$ ambient coordinates,
with no separate degree restriction. Unused coordinates and boundary
points are permitted. Fixed coordinates ($\ell_i=u_i$) are permitted in
both definitions. Constants can be handled without division by a zero gap.

The paper proceeds from an elementary disproof to the exact finite gap
and then to the sharp replacement. Appendix~\ref{app:coupling} gives the
complete common-coupling upper argument. Appendix~\ref{app:lean} connects
the statements to the bundled Lean declarations. All mathematical
ingredients needed for the argument are proved below; the bibliography
identifies the conjecture and its immediate context.

\section{Graph hulls and probability laws}\label{sec:foundations}
We omit the subscript $B$ on a unit cube. A polynomial is multilinear
when it is affine in each coordinate separately.

\begin{lemma}[Vertex laws]\label{lem:vertices}
For a multilinear $g$ on a finite box $B$, let $\laws_B(x)$ be the
probability laws on its vertices having mean $x$. Then
\[
 \vex_B g(x)=\min_{P\in\laws_B(x)}\E_Pg(X),\qquad
 \cav_B g(x)=\max_{P\in\laws_B(x)}\E_Pg(X).
\]
Both extrema are attained, and these are the endpoints of the vertical
slice of the original continuous graph hull.
\end{lemma}
\begin{proof}
Round each nonfixed coordinate of a point $z\in B$ independently to its
two endpoints, preserving its mean. Repeated use of separate affinity
gives $\E X=z$ and $\E g(X)=g(z)$. A fixed coordinate stays fixed.
Thus every continuous graph point lies in the convex hull of the vertex
graph. The reverse hull inclusion is immediate. A convex combination of
vertex graph points is exactly a vertex law with its mean and expected
function value. Laws with a specified mean form a nonempty compact
polytope; hence minimizing and maximizing the expected value attain the
two endpoints.
\end{proof}

\begin{lemma}[Individual envelopes and a common upper law]\label{lem:monomial}
For a nonempty support $e$ on the unit cube,
\begin{equation}\label{eq:monomial}
 \cav m_e(x)=\min_{i\in e}x_i,
 \qquad \vex m_e(x)=\max\left(0,\sum_{i\in e}x_i-|e|+1\right).
\end{equation}
A single law attains the upper envelope of every monomial simultaneously.
Consequently, $\cav f(x)=\sum_e a_e\cav m_e(x)$ when $a_e\ge0$.
\end{lemma}
\begin{proof}
On a binary vertex, a monomial indicates simultaneous success of all its
coordinates. The joint-success probability is at most the smallest
success marginal and at least one minus the sum of the failure marginals,
or zero if that is larger. For the upper bound, take one uniform
$U\in[0,1]$ and set $X_i=\ind\{U\le x_i\}$ for every coordinate.
This attains all upper bounds at once.

For the lower bound for one support, place intervals of lengths $1-x_i$
consecutively on a circle of circumference one. Their union has length
$\min(1,\sum_{i\in e}(1-x_i))$: before one full turn the intervals are
disjoint up to endpoints, and after a full turn they cover the circle.
Let each interval be the corresponding failure event. This realizes
the claimed joint-success probability. Sample unused coordinates with
their prescribed means to complete the law. Nonnegative coefficients
allow the common upper equalities to be summed. A constant term is
immediate.
\end{proof}

These elementary envelope facts are also part of the classical
multilinear relaxation theory discussed in \citet{Luedtke2012}. The
essential distinction is that each term can have its own minimizing law
in $T$, while the convex envelope of the sum requires one law for all
terms together.

\section{A sparse family with unbounded ratio}\label{sec:disproof}
Fix an integer $L\ge2$ and put $m=2^L$. Use leaves $z_0,\ldots,z_{m-1}$
and anchors $a_1,\ldots,a_L$. At level $j$, partition the leaf indices
into consecutive blocks
\[
 \mathcal B_j=\bigl\{\{b2^{L-j},\ldots,(b+1)2^{L-j}-1\}:
                       0\le b<2^j\bigr\}.
\]
The partitions are nested. Define the polynomial and evaluation point by
\begin{equation}\label{eq:family}
 f_L(a,z)=\sum_{j=1}^L\sum_{B\in\mathcal B_j}a_j\prod_{i\in B}z_i,
 \qquad a_j=2^{-j},\quad z_i=1-2^{-L}.
\end{equation}
We write $T_L$ and $H_L$ for its two gaps at this point. Each support is
distinguished by its anchor and its block, so all coefficients are one
on distinct squarefree supports. The number of variables, number of
monomials, maximum degree, and number of variable occurrences are
\begin{equation}\label{eq:sizes}
 n_L=2^L+L,\quad |E_L|=2^{L+1}-2,\quad
 d_L=2^{L-1}+1,\quad
 \sum_{e\in E_L}|e|=L2^L+2^{L+1}-2.
\end{equation}
The counts follow by summing $2^j$ terms of size $2^{L-j}+1$ at each
level. Every prescribed coordinate lies strictly between zero and one.

\begin{lemma}[Termwise gap and failure representation]\label{lem:family-gap}
For the family \eqref{eq:family}, $T_L=L$, the upper envelope is $L$,
and $H_L>0$. Under any binary law with the prescribed means, let $A_j$
be the anchor, $R$ the number of failed leaves, and $N_j$ the number of
level-$j$ blocks containing a failed leaf. Then
\begin{equation}\label{eq:failure-representation}
 \E R=1,\quad \E A_j=2^{-j},\quad
 H_L=\max_P\E_P\sum_{j=1}^L A_jN_j,
 \qquad N_j\le\min(2^j,R).
\end{equation}
The maximum ranges over all binary laws with all original coordinate
means, without independence assumptions.
\end{lemma}
\begin{proof}
At level $j$, a support has $k_j=2^{L-j}$ leaves and one anchor. Its
lower envelope is $\max(0,2^{-j}-k_j2^{-L})=0$, and its upper envelope
is $2^{-j}$. There are $2^j$ such terms. Lemma~\ref{lem:monomial}
therefore gives $T_L=\cav f_L=L$. The value at the original point is
\[
 f_L(a,z)=\sum_{j=1}^L(1-2^{-L})^{2^{L-j}}<L.
\]
Since $\vex f_L\le f_L$, the hull gap is positive.

The expected number of failures is $m2^{-L}=1$. At every binary vertex,
$f_L=\sum_j A_j(2^j-N_j)$. Taking expectations gives
$\E f_L=L-\E\sum_j A_jN_j$. Apply Lemma~\ref{lem:vertices} to minimize
$\E f_L$. Finally, the blocks of each partition are disjoint, so each
hit block needs a different failed leaf.
\end{proof}

The next bound applies directly to every law in
\eqref{eq:failure-representation}. It requires only a finite sum.

\begin{lemma}[A finite cap certificate]\label{lem:cap}
For $r\ge0$, $A_j\in\{0,1\}$, and integers $0\le s\le L$,
\begin{equation}\label{eq:cap}
 \sum_{j=1}^L A_j\min(2^j,r)
 \le sr+\sum_{j=s+1}^L A_j2^{j-s}.
\end{equation}
Consequently,
\begin{equation}\label{eq:finite-hull-bound}
 H_L\le s+(L-s)2^{-s}\qquad(0\le s\le L).
\end{equation}
\end{lemma}
\begin{proof}
Put $c_j=\min(2^j,r)$. For $j>s$,
\[
 A_jc_j\le c_j-c_{j-s}+A_j2^{j-s}.
\]
If $A_j=0$ this uses $c_j\ge c_{j-s}$; if $A_j=1$ it uses
$c_{j-s}\le2^{j-s}$. Bound the first $s$ selected caps by $c_j$ and
sum. The cap differences telescope to
$\sum_{j=L-s+1}^L c_j\le sr$, proving \eqref{eq:cap}.
Now use $N_j\le\min(2^j,R)$ and take expectations. Each term in the
last sum has expectation $2^{j-s}2^{-j}=2^{-s}$, and $\E R=1$.
\end{proof}

\begin{theorem}[Disproof of the uniform-constant conjecture]\label{thm:disproof}
For every real $K$ there is a unit-coefficient multilinear polynomial on
a unit cube, at a strictly interior point with positive hull gap, for
which $T/H>K$. In particular, the Luedtke--Namazifar--Linderoth
uniform-constant conjecture is false.
\end{theorem}
\begin{proof}
Choose $s=\lceil\log_2 L\rceil$ in \eqref{eq:finite-hull-bound};
then $s\le L$ and $(L-s)2^{-s}\le1$. Hence
\[
 \frac{T_L}{H_L}\ge\frac{L}{\lceil\log_2 L\rceil+1}\longrightarrow\infty.
\]
All coefficient, support, interiority, and positive-denominator
conditions were checked in \eqref{eq:family} and
Lemma~\ref{lem:family-gap}. A unit cube is an admissible nonnegative box.
\end{proof}

The degree grows with $L$. This theorem is consistent with finite
degree-dependent bounds, whose sharp growth is determined below. A
single example exceeding a particular numerical ratio would not suffice
to disprove the existence of every universal constant; the family does.

\section{The exact hull gap and an attaining law}\label{sec:exact}
The finite bound is attained at a suitable cutoff. Define
\[
 B_q=(L-q+2)2^{-q}\qquad(1\le q\le L).
\]
This sequence decreases strictly, with $B_1\ge1\ge B_L$.
There is therefore an integer $1\le s<L$ such that
\begin{equation}\label{eq:cutoff}
 B_{s+1}\le1\le B_s,
 \quad\text{or equivalently}\quad
 (L-s+1)2^{-(s+1)}\le1\le(L-s+2)2^{-s}.
\end{equation}

\begin{theorem}[Exact finite gap]\label{thm:exact}
For every $L\ge2$ and every cutoff $s$ satisfying \eqref{eq:cutoff},
\begin{equation}\label{eq:exact}
 T_L=L,\qquad H_L=s+(L-s)2^{-s}.
\end{equation}
Both envelope endpoints are attained. Either valid adjacent cutoff at
an equality in \eqref{eq:cutoff} gives the same value.
\end{theorem}
\begin{proof}
The upper bound for $H_L$ is Lemma~\ref{lem:cap}. To attain it, first
choose $l\in\{0,\ldots,L\}$ with probabilities
\begin{equation}\label{eq:threshold-weights}
 w_l=\begin{cases}2^{-l-1},&l<L,\\2^{-L},&l=L.\end{cases}
\end{equation}
Set $A_j=\ind\{j\le l\}$. Summing the geometric tail gives
$\Prb(A_j=1)=2^{-j}$. For a cutoff $q\in\{s,s+1\}$, prescribe
\begin{equation}\label{eq:failure-cutoff}
 R_l^{(q)}=\begin{cases}0,&l<q,\\2^{l-q+1},&l\ge q.\end{cases}
\end{equation}
The failure count is at most $2^L$, since $q\ge1$. A direct sum gives
\[
 \E R^{(q)}=(L-q)2^{-q}+2^{1-q}=B_q.
\]
Independently of $l$, mix the two choices of $q$ with weights
\begin{equation}\label{eq:cutoff-mix}
 \theta=\frac{1-B_{s+1}}{B_s-B_{s+1}},\qquad
 1-\theta=\frac{B_s-1}{B_s-B_{s+1}}.
\end{equation}
They lie in $[0,1]$ and make the mixed failure count have mean one.

We next realize the block hits and each individual leaf mean. If the
chosen failure count is zero, let every leaf succeed. If it is
$R=2^r$, choose a uniform residue $t$ modulo $2^{L-r}$ and fail exactly
the leaves whose indices are congruent to $t$ modulo $2^{L-r}$.
This gives $2^r$ failed leaves and conditional failure probability
$R/2^L$ for every leaf. At level $j$, a block has length $2^{L-j}$.
If $j\le r$, each block contains a failure, so the number hit is $2^j$.
If $j>r$, the spacing between failures exceeds the block length, and
the $2^r$ failures hit different blocks. In both cases,
\begin{equation}\label{eq:all-hits}
 N_j=\min(2^j,R)
\end{equation}
simultaneously for all levels. Averaging $R/2^L$ over the mixed law
gives failure probability $2^{-L}$ for each leaf. Thus the entire law
has the original prescribed means.

It remains to check the attained objective. For the nested anchors put
$S_l(R)=\sum_{j=1}^l\min(2^j,R)$. For $l>s$ and either
$R=2^{l-s}$ or $R=2^{l-s+1}$,
\[
 S_l(R)=sR+2^{l-s+1}-2
       =sR+\sum_{j=s+1}^l2^{j-s}.
\]
Indeed, cap the first $l-s$ terms at $2^j$ and the last $s$ at $R$;
the two endpoint choices both give equality. These are exactly the
failure counts in \eqref{eq:failure-cutoff}. For $l=s$ they are $0$
and $2$, both satisfying $S_s(R)=sR$; for $l<s$ they are both zero.
Therefore \eqref{eq:cap} is an equality in every state of the mixed law.
Taking expectations and using \eqref{eq:all-hits} gives
$\E\sum_j A_jN_j=s+(L-s)2^{-s}$. This attains the maximum in
\eqref{eq:failure-representation}, and hence attains the lower envelope
$L-H_L$. The common upper law attains $L$.

Finally, the formulas at $s$ and $s+1$ differ by $B_{s+1}-1$, so they
agree whenever both adjacent choices are valid.
\end{proof}

\begin{table}[htbp]
\centering
\begin{tabular}{rrrr}
\toprule
$L$ & A valid $s$ & $H_L$ & $T_L/H_L$\\
\midrule
2 & 1 & $3/2$ & $4/3$\\
8 & 2 & $7/2$ & $16/7$\\
16 & 3 & $37/8$ & $128/37$\\
64 & 5 & $219/32$ & $2048/219$\\
\bottomrule
\end{tabular}
\caption{Exact instances of Theorem~\ref{thm:exact}. These values
illustrate the formula; the universally quantified proof establishes
unboundedness.}\label{tab:examples}
\end{table}

The cutoff also explains the growth. From $2^s\le L-s+2\le L+1$,
$s\le\log_2(L+1)$. The other cutoff inequality gives
$2^{s+1}\ge L-s+1$, which, together with this upper bound, yields
$s=\log_2 L+O(1)$. Moreover,
\[
 0\le (L-s)2^{-s}\le 2\frac{L-s}{L-s+1}<2.
\]
Consequently,
\begin{equation}\label{eq:family-asymptotic}
 H_L=\log_2 L+O(1),\qquad
 \frac{T_L}{H_L}\sim\frac{L}{\log_2 L}.
\end{equation}
Nesting of the partitions was used for simultaneous attainment in
\eqref{eq:all-hits}. The finite bound and disproof themselves only need
the specified equal block sizes and disjointness within each level.

\section{The sharp replacement}\label{sec:sharp}
The construction gives the sharp leading scale for both degree and
dimension. Here and below $\ln$ denotes the natural logarithm.

\begin{theorem}[Sharp degree and dimension growth]\label{thm:sharp}
For the suprema defined in Section~1,
\begin{equation}\label{eq:sharp}
 R(d)\sim\frac{\ln d}{\ln\ln d},\qquad
 C(n)\sim\frac{\ln n}{\ln\ln n}.
\end{equation}
Both limits hold along all positive integer allowances tending to
infinity, with leading constant one, over all finite nonnegative boxes.
\end{theorem}

We state the explicit bound used for the upper half. It matches the
finite harmonic certificate in the accompanying Lean proof.

\begin{theorem}[A common-coupling finite bound]\label{thm:finite-upper}
Let $d\ge2$, $c_0=1-e^{-1}$, and define
\begin{align}
 \Lambda&=\max\{e^6,1+\ln d\},& b&=\ln\Lambda,& A&=\Lambda/b^2,
       \label{eq:upper-parameters}\\
 h&=\frac{1-1/b}{1+1/b^2}\,
          \frac{b-3\ln b}{\Lambda},&
 Z&=\frac1h+A+\frac2{c_0}.\label{eq:upper-constant}
\end{align}
Then $h>0$ and $T_B(f;x)\le ZH_B(f;x)$ for every positive multilinear
polynomial of degree at most $d$ on a finite nonnegative box, including
points of zero hull gap. On a unit cube the bound is witnessed by a
single law depending only on the complete marginal vector and $d$,
independently of the supports and coefficients.
\end{theorem}

Appendix~\ref{app:coupling} proves this theorem. The finite constant is
chosen to give the sharp leading term; no optimality of $Z$ at a fixed
degree is asserted.

\begin{proof}[Proof of Theorem~\ref{thm:sharp}]
For the upper bound, $b\to\infty$ and
\[
 \frac{h}{b/\Lambda}
   =\frac{1-1/b}{1+1/b^2}\left(1-\frac{3\ln b}{b}\right)
      \longrightarrow1.
\]
Thus $1/h\sim\Lambda/\ln\Lambda$, while
$A/(\Lambda/\ln\Lambda)=1/b\to0$ and the constant $2/c_0$ is
negligible. Eventually $\Lambda=1+\ln d$, so Theorem~\ref{thm:finite-upper}
gives
\[
 R(d)\le(1+o(1))\frac{\ln d}{\ln\ln d}.
\]
Every support in $n$ coordinates has size at most $n$, so $C(n)\le R(n)$.

For a common lower argument, take an integer $N\ge8$ and
$L=\lfloor\log_2N\rfloor-1$. Then $L\ge2$ and
$n_L=2^L+L\le2^{L+1}\le N$. The family therefore fits degree allowance
$N$. It also fits exactly $N$ coordinates after adding unused variables
with arbitrary fixed cube means, for instance $1/2$. Projection of any
law on the added coordinates gives an old law, and any old law extends
by independent added coordinates. Both gaps are therefore unchanged.
By \eqref{eq:family-asymptotic}, both $R(N)$ and $C(N)$ are at least
\[
 \frac{T_L}{H_L}\sim\frac{L}{\log_2L}
     \sim\frac{\ln N}{\ln\ln N}.
\]
This construction works for every $N\ge8$. Together with the upper
bounds it proves both limits. In particular, the relevant ratio sets
are nonempty and bounded above before taking their suprema.
\end{proof}

The conjectured constant is therefore replaced by a slowly growing
degree- or dimension-dependent factor. The lower family has fewer than
$2n_L$ terms and $O(n_L\log n_L)$ variable occurrences. These elementary
size consequences do not claim an efficient algorithm for constructing
or evaluating a general convex envelope.

\appendix
\section{The common-coupling upper bound}\label{app:coupling}
We give all details needed for Theorem~\ref{thm:finite-upper}. The
parameters and the rational union estimate below follow the formal
proof, so no optimized fixed-point or second-order certificate is needed.

\subsection{Deficiencies and the easy terms}
Discard constant and affine terms, which have zero width and do not
change $H$. For a remaining support $e$, choose an anchor $i_e$ of
minimum mean and write
\begin{equation}\label{eq:deficiency}
 u_e=x_{i_e},\quad p_j=1-x_j\ (j\ne i_e),\quad
 t_e=\min\left(u_e,\sum_{j\in e\setminus\{i_e\}}p_j\right),\quad
 D_e(X)=X_{i_e}-\prod_{i\in e}X_i.
\end{equation}
For binary $X$, $D_e$ is the nonnegative indicator that the anchor
succeeds and at least one other coordinate fails. Lemmas~\ref{lem:vertices}
and~\ref{lem:monomial} imply
\begin{equation}\label{eq:deficiency-gap}
 T(f;x)=\sum_e a_et_e,\qquad
 H(f;x)=\max_{P\in\laws(x)}\sum_e a_e\E_PD_e.
\end{equation}
Thus a law satisfying $\E D_e\ge t_e/Z$ for every support bounds the
ratio after summing with nonnegative coefficients. The law must be
chosen once for the entire vector $x$.

Call a coordinate low if $x_i\le1/2$ and high otherwise. Independent
Bernoulli rounding gives
\[
 \E D_e=u_e\left(1-\prod_{j\in e\setminus\{i_e\}}(1-p_j)\right).
\]
If there are at least two low coordinates in the support, a nonanchor
has $p_j\ge1/2$, so the deficiency is at least $u_e/2\ge t_e/2$.
If there are no low coordinates, let $S=\sum_jp_j$. The elementary
inequality $1-p\le e^{-p}$ gives
\[
 \E D_e\ge u_e(1-e^{-S})\ge c_0u_e\min(1,S)\ge c_0t_e/2.
\]
For the middle inequality, use concavity of $1-e^{-S}$ on $[0,1]$
and monotonicity after $1$. For the last, if $S\le1$ use $u_e>1/2$;
otherwise use $t_e\le u_e$. Independence therefore captures at least
$c_0t_e/2$ in these two cases. Terms of zero gap need no division and
satisfy the desired inequality under every law.

\subsection{Two laws for the one-low case}
Put $M=e^{\Lambda-1}\ge d$, so $1+\ln M=\Lambda$. Each of the next two
laws draws one uniform $U\in(0,1)$ and sets every low coordinate to
$X_i=\ind\{U\le x_i\}$.

In the threshold law, a high coordinate with failure mean $p=1-x_i$
fails exactly when $U\le p$. All means are preserved. For a one-low
support with low-anchor mean $u$, this law gives deficiency
\begin{equation}\label{eq:threshold-gain}
 \min\left(u,\max_jp_j\right).
\end{equation}

In the harmonic law, a high coordinate with $p=0$ never fails. For
$p>0$, put
\begin{equation}\label{eq:harmonic-density}
 h_p=\min(Mp,1),\quad a_p=1+\ln(h_p/p),\quad
 q_p(t)=\frac{\min(1,p/t)}{a_p}\ind\{t\le h_p\}
       \quad(0<t<1).
\end{equation}
Conditional on $U=t$, high coordinates fail independently with
probabilities $q_p(t)$. Since $p\le h_p$ and $1\le a_p\le\Lambda$,
these probabilities belong to $[0,1]$. They are measurable and
\begin{equation}\label{eq:density-integral}
 \int_0^1q_p(t)\,dt
 =\frac{p+p\ln(h_p/p)}{a_p}=p.
\end{equation}
Hence every high coordinate also has the prescribed mean. Values at
the endpoints of the integral have no effect. The finite binary law
can equivalently be defined by integrating the conditional probability
of each binary vertex; finite sums commute with the integral.

For a one-low support its anchor is the low coordinate. Conditional
on $U=t\le u$, the probability of a nonanchor failure is
$1-\prod_j(1-q_{p_j}(t))$. Its harmonic deficiency is therefore
\begin{equation}\label{eq:harmonic-deficiency}
 \int_0^u\left[1-\prod_j(1-q_{p_j}(t))\right]dt.
\end{equation}
Both laws are defined using the complete marginal vector, without
reference to a support or its coefficient.

\subsection{The harmonic window and the mixture}
First, for independent events of probabilities $q_j\in[0,1]$ and
$S=\sum_jq_j$,
\begin{equation}\label{eq:rational-union}
 1-\prod_j(1-q_j)\ge\frac{S}{1+S}.
\end{equation}
Indeed, $\prod_j(1+q_j)\ge1+S$ and
$\prod_j(1-q_j)\prod_j(1+q_j)=\prod_j(1-q_j^2)\le1$.
Multiply the first inequality by the nonnegative failure-complement
product and rearrange. The right side is increasing in $S\ge0$.

Let $\beta=1/b$, with $b,A,h$ as in
\eqref{eq:upper-parameters}--\eqref{eq:upper-constant}. We have $b\ge6$,
and $b-3\ln b>0$: it is positive at $6$ (since $e^2>6$) and its
derivative is $1-3/b>0$ thereafter. Thus $h>0$ and
\[
 \ln(\beta A)=b-3\ln b>0,\qquad \beta A>1,\qquad A>1.
\]
For a one-low term with $t_e>0$, if some $p_j\ge t_e/A$, then
\eqref{eq:threshold-gain} is at least $t_e/A$, since $t_e\le u$.
Otherwise every $p_j<t_e/A$. Consider the nonempty window
\begin{equation}\label{eq:window}
 t_e/A\le t\le\beta t_e\le u.
\end{equation}
All $p_j\le t$ there. Call a coordinate active when $t\le Mp_j$.
Since $t\le1$, this is the same as $t\le h_{p_j}$. Every inactive
coordinate has $p_j<t/M$, and there are at most $d-1\le M$ coordinates
in the nonanchor support. The total inactive failure mass is at most
$t$, so the active mass is at least
\[
 \sum_jp_j-t\ge t_e-t\ge(1-\beta)t_e.
\]
On active coordinates, \eqref{eq:harmonic-density} gives
$q_{p_j}(t)\ge p_j/(\Lambda t)$. Consequently,
\[
 \sum_jq_{p_j}(t)\ge z(t):=\frac{(1-\beta)t_e}{\Lambda t},
 \qquad 0\le z(t)\le\frac A\Lambda=\frac1{b^2}.
\]
Using \eqref{eq:rational-union}, the conditional union probability is
at least $z(t)/(1+z(t))\ge z(t)/(1+1/b^2)$. Integrate over
\eqref{eq:window}, using nonnegativity outside that window, to obtain
\begin{align*}
 \E_{\mathrm{harm}}D_e
 &\ge\frac{1-\beta}{1+1/b^2}\frac{t_e}{\Lambda}
              \ln(\beta A)\\
 &=\frac{1-1/b}{1+1/b^2}\frac{b-3\ln b}{\Lambda}\,t_e
   =ht_e.
\end{align*}
The integrands are bounded measurable functions on a finite interval,
and the reciprocal integration window stays strictly above zero.

Mix the harmonic, threshold, and independent laws with weights
\begin{equation}\label{eq:three-mixture}
 \frac{1/h}{Z},\qquad\frac A Z,\qquad\frac{2/c_0}{Z},
 \qquad Z=1/h+A+2/c_0.
\end{equation}
These positive weights sum to one and preserve every coordinate mean.
Each positive-gap term receives at least $t_e/Z$ from one of the three
components. All other components contribute nonnegative deficiencies.
Summing \eqref{eq:deficiency-gap} proves $T\le ZH$ on the cube,
including at zero hull gap. This proves the cube part of
Theorem~\ref{thm:finite-upper}.

\subsection{Transfer to original terms on a nonnegative box}
For each coordinate substitute
$x_i=\ell_i+(u_i-\ell_i)y_i$ with $0\le y_i\le1$, and let
$\widetilde f(y)=f(x)$. Expanding original monomials gives nonnegative
coefficients and does not increase degree.

First suppose all coordinates are nonfixed. The substitution is an
affine bijection, so it maps graph hulls and preserves their vertical
widths:
\begin{equation}\label{eq:box-transfer}
 H_B(f;x)=H(\widetilde f;y),\qquad
 T_B(f;x)\le T(\widetilde f;y).
\end{equation}
For the second inequality, an original term becomes a sum of expanded
monomials. The sum of their convex envelopes is a convex underestimator
of that term; the sum of their concave envelopes is a concave
overestimator. Thus the original term's width is at most the sum of
the expanded widths. Sum this inequality over the original terms.
If expansion produces the same support more than once, combining its
nonnegative coefficients leaves that sum of widths unchanged, because
envelopes scale linearly by a nonnegative constant.

If a coordinate is fixed, restrict the original box to its nonfixed
coordinates. In the expanded function the fixed coordinate has zero
width and its new cube coordinate is unused. Projecting or extending
laws on unused coordinates preserves every attainable expectation, as
in the proof of Theorem~\ref{thm:sharp}. Thus both the full-hull equality
and the original-term inequality in \eqref{eq:box-transfer} still hold.
If all coordinates are fixed, every gap is zero. Applying the cube
bound to $\widetilde f$ now proves
$T_B(f;x)\le T(\widetilde f;y)\le ZH(\widetilde f;y)=ZH_B(f;x)$,
completing Theorem~\ref{thm:finite-upper}.

\section{Formal statements and reproduction}\label{app:lean}
The accompanying \href{formal/README.md}{\texttt{formal/}} directory
is a standalone Lake project pinned to Lean~4.33.1 and
Mathlib~v4.33.1, with exact dependency revisions in
\texttt{lake-manifest.json}. No parent repository files are needed.
The 48 proof modules consist of 41 multilinear modules and seven shared
modules for finite laws, graph hulls, and individual envelopes. The
shared namespace \texttt{CubicGap} is inherited from their original
location; its bundled files are general infrastructure, and no cubic
counterexample modules are included.

The following source links are relative to this PDF. They work when the
PDF and the accompanying folder are kept together and the viewer permits
local links. File paths and declaration names also appear in
\href{formal/COVERAGE.md}{\texttt{formal/COVERAGE.md}} for viewers that
disable such links. No hosted verification service is required.

\begin{longtable}{@{}>{\raggedright\arraybackslash}p{.30\textwidth}>{\raggedright\arraybackslash}p{.65\textwidth}@{}}
\toprule
Paper claim & Lean file and principal declarations\\
\midrule
\endhead
Lemma~\ref{lem:vertices}, continuous cube hull &
\leanfile{CubicGap/Hull}{Hull.lean}: \texttt{mem\_cubeGraph\_hull\_iff}.\\[4pt]
Lemma~\ref{lem:family-gap}, $T_L=L$ and $H_L>0$ &
\leanfile{MultilinearGap/Termwise}{Termwise.lean}: \texttt{termwiseGap\_eq};
\leanfile{MultilinearGap/Results}{Results.lean}: \texttt{hullGap\_positive}.\\[4pt]
Lemma~\ref{lem:cap} &
\leanfile{MultilinearGap/ExactUpper}{ExactUpper.lean}:
\texttt{dyadic\_exact\_pointwise}, \texttt{hullGap\_exact\_upper\_bound}.\\[4pt]
Theorem~\ref{thm:disproof} &
\leanfile{MultilinearGap/Results}{Results.lean}:
\texttt{unbounded\_gap\_ratio},
\texttt{no\_uniform\_positive\_multilinear\_bound}.\\[4pt]
Theorem~\ref{thm:exact} &
\leanfile{MultilinearGap/ExactResults}{ExactResults.lean}:
\texttt{polynomial\_exact\_minimum},
\texttt{hullGap\_exact}, \texttt{exists\_hullGap\_exact}.\\[4pt]
Theorem~\ref{thm:finite-upper}, cube &
\leanfile{MultilinearGap/SharpUpper}{SharpUpper.lean}:
\texttt{sharpLaw\_hasMeans}, \texttt{sharp\_cube\_gap\_bound}.\\[4pt]
Theorem~\ref{thm:finite-upper}, original boxes &
\leanfile{MultilinearGap/BoxTransfer}{BoxTransfer.lean}:
\texttt{degree\_gap\_bound\_on\_box}.\\[4pt]
Theorem~\ref{thm:sharp} &
\leanfile{MultilinearGap/SharpAsymptotics}{SharpAsymptotics.lean}:
\texttt{sharp\_positive\_growth}.\\
\bottomrule
\end{longtable}

All multilinear declarations in the table belong to namespace
\texttt{MultilinearGap}. The shared hull declarations belong to
\texttt{CubicGap}. The final asymptotic theorem states the two normalized
limits for the actual all-box degree and exactly-$n$-coordinate suprema.
It is not a theorem merely along the dyadic family.

The formal disproof endpoint uses the elementary weaker estimate
$H_L\le s+2L/2^s$ and an arithmetic parameter choice, independently of
the exact formula. This paper uses the sharper finite certificate,
which is also verified, to shorten the written disproof. Similarly,
the formal asymptotic lower argument uses the verified estimate
$H_L\le\log_2 L+3$ and does not depend on the exact attaining law.
The conclusions agree; these are alternative proof routes within the
same bundle.

The completion modules also verify the paper's supporting mathematical
claims.

\leanfile{MultilinearGap/Attainment}{Attainment.lean} proves general
endpoint attainment, and
\leanfile{MultilinearGap/EnvelopeFunctions}{EnvelopeFunctions.lean}
identifies those endpoints with the greatest convex underestimator and
least concave overestimator on both cubes and boxes.
\leanfile{MultilinearGap/MonomialEnvelope}{MonomialEnvelope.lean}
proves the full individual lower-envelope identity in
Lemma~\ref{lem:monomial}, including an attaining law preserving all ambient
means.

The module \leanfile{MultilinearGap/PaperFoundations}{PaperFoundations.lean}
checks original-box coefficient scaling and $0\le H_B\le T_B$.

\leanfile{MultilinearGap/ExactSize}{ExactSize.lean} proves all four size
identities in \eqref{eq:sizes}, with the coordinate count inherited from
\texttt{FamilySize.lean}, and the stated sparsity bounds.
\leanfile{MultilinearGap/ExactAsymptotics}{ExactAsymptotics.lean} proves
$|H_L-\log_2L|\le3$ for $L\ge2$ and the actual family ratio limit in
\eqref{eq:family-asymptotic}.
\leanfile{MultilinearGap/Examples}{Examples.lean} proves all four rows of
Table~\ref{tab:examples} as equalities about the actual graph hull.

With Elan, Git, and Python~3 installed and the pinned toolchain on
\texttt{PATH}, run from \texttt{formal/}:
\begin{verbatim}
lake exe cache get
bash scripts/verify.sh
\end{verbatim}
The first command obtains the pinned dependency cache. The script checks
that every proof module is imported, builds with warnings treated as
failures, audits all bundled declarations (including private helpers)
for transitive axiom dependencies, and replays the compiled project
declarations through Lean's kernel. Only \texttt{propext},
\texttt{Classical.choice}, and \texttt{Quot.sound} are permitted.
Unfinished proofs, custom axioms, and native-computation axioms are
rejected. No external optimization solver or sampled numerical result
is a trusted proof step.

The replay uses the installed Lean kernel and compiled dependencies;
it is not a separate kernel implementation and does not replay all
of Mathlib from scratch. The mathematical interpretation also depends
on matching the formal definitions to the statements above. The
\href{formal/VERIFICATION.md}{verification record} specifies the actual
local checks and their scope. Internal agent reviews of the antecedent
proofs do not constitute external peer review or establish publication
priority. This note makes no exhaustive priority claim.

The formal scope is the disproof, exact family gap and attainment, the
finite bound \eqref{eq:upper-constant}, and the sharp leading limits.
It does not include optimized fixed-point or Lambert $W$ formulas,
second-order asymptotics, coefficient removal, or homogenization.
Those results are not needed for any theorem stated here.

\bibliographystyle{plainnat}
\bibliography{references}
\end{document}
